\section{The critical surface exponent}
\label{sec:v89-critical-surface}
Put $\gamma=1/32$ and $Q=(1-\gamma)^{-1}=32/31$. We fix
$\eta=1/4$, $\kappa=3/8$ and $d=1-\kappa=5/8$ in this section.
The endpoint estimate does not integrate a negative residence moment
of order $\gamma$. Instead it retains the size of the tested roof set
until after summing the complete residence distribution.

\begin{lemma}[A capped dyadic pulse sum]\label{lem:v89-pulse-sum}
Let $s_j=2^{-j}s$, $0<\theta<1$ and $v,s>0$. Then
\begin{equation}\label{eq:v89-pulse-sum}
 \sum_{j\ge0}\min\{1,2v/s_j\}s_j^\theta
                 \le C_\theta\min\{v^\theta,s^\theta\}.
\end{equation}
For $0<\gamma<a<1$ and $0<D\le1$, one also has
\begin{equation}\label{eq:v89-capped-sum}
 \sum_{j\ge0}\min\{1,2v/s_j\}\min\{D,s_j^a\}
                 \le C_\gamma D^{1-\gamma/a}v^\gamma.
\end{equation}
All sums are formed before taking a supremum over $v$.
\end{lemma}
\begin{proof}
For $v<s$, split at the first index for which $s_j\le v$.
The later terms sum a geometric series of order $v^\theta$.
The earlier terms are bounded by $2v s_j^{\theta-1}$; their
geometric sum has the same order. The full sum is also bounded by
$\sum_j s_j^\theta=s^\theta/(1-2^{-\theta})$. This proves
\eqref{eq:v89-pulse-sum}. Finally
$\min(D,x^a)\le D^{1-\gamma/a}x^\gamma$ for every $x>0$.
Apply the first assertion at exponent $\gamma$.
\end{proof}

Write $Z_s$ for the complete outgoing pulse source defined in
Theorem~\ref{thm:v88-weighted-trace}. Its traces are the actual
one-sided nongrazing outgoing traces, including every incoming type.
Let $J^{[s]}=J+(-s,s)$, $|J|=L$, and let
$V_m(J^{[s]})=m^2\int_{J^{[s]}}v^{1/4}(u)\,du$.

\begin{theorem}[Finite-count critical distribution bound]
\label{thm:v89-endpoint-surface}
For every fixed auxiliary band $B_*$, all $m\ge2$ and
$0<s<s_0$, the full short surface satisfies
\begin{align}
 \mathcal N_Q(m^2Z_s;J)\le{}&C(L+2s+B_*^{-1})
           +Cs^{19/32}V_m(J^{[s]})\notag\\
 &+C_{B_*}(1+B_*(L+2s))m^3\rho_{B_*}^{15m/32}.
                                      \label{eq:v89-critical-finite}
\end{align}
The bound holds uniformly over exact labels, radius and interval
location. It includes the supremum over every measurable subset of
$J$, even one depending on $m$, before either limit. Consequently
\begin{equation}\label{eq:v89-critical-ordered}
 \limsup_m\sup_{R,n,k}\sup_{|J|=L}\mathcal N_Q(m^2Z_s;J)
                 \le C(L+2s)(1+s^{19/32}).
\end{equation}
There is no positive lower residence, incidence or gradient cutoff.
\end{theorem}
\begin{proof}
For one maximal segment $\sigma$, put
$T_\sigma=A_\sigma(b_\sigma-)$,
$U_\sigma=\int_{a_\sigma}^{b_\sigma}A_\sigma$ and
$V_\sigma=\int_{a_\sigma}^{b_\sigma}A_\sigma I_m z^{-1/4}$.
The whole-segment calculation in
\eqref{eq:v88-trace-triangle}, without a residence weight, gives
\begin{equation}\label{eq:v89-unweighted-trace}
             T_\sigma\ell_\sigma
                   \le U_\sigma+C\ell_\sigma^d V_\sigma.
\end{equation}
It differentiates the complete density $\omega/(4\pi c_*\chi)$.
The factor $u-a_\sigma$ obtained by integration over starting roofs
is what permits the radial lower bound to be used at a critical
incoming end. No inverse physical Jacobian is discarded.

Fix a measurable $E\subset J$ of measure $v>0$. A segment in the
band $s_j/2<\ell_\sigma\le s_j$ contributes at most
$T_\sigma\min(v,\ell_\sigma)$ to $\int_E Z_s$.
Its contribution is therefore at most
\[
 \min\{1,2v/s_j\}\{U_\sigma+Cs_j^dV_\sigma\}.
\]
Every such segment intersecting $J$ lies in $J^{[s]}$.
The Borel segment bijection counts its interior once. Thus the sum
of $U_\sigma$ in that band is at most the original volume mass
$\bar\nu\{F\in J^{[s]},\ell\le s_j\}$, and the sum of
$V_\sigma$ is at most $\int_{J^{[s]}}v^{1/4}$.
These are positive upper comparisons on the same exact event;
no occupation coefficients are combined.

Insert \eqref{eq:v88-capped-residence} in the first sum. Its leading
part is bounded by $C(L+2s+B_*^{-1})v^\gamma$, by
\eqref{eq:v89-pulse-sum}. In its remainder apply
\eqref{eq:v89-capped-sum} with $a=1/2$ and
$D=\rho_{B_*}^{m/2}$. Since $1-2\gamma=15/16$, its bound is
\[
 C_{B_*}(1+B_*(L+2s))m^3\rho_{B_*}^{15m/32}v^\gamma.
\]
This is the complete finite-count tail, not a sum of separate
collision limsups. For the variation part use
\eqref{eq:v89-pulse-sum} at exponent $d$ and
$\min(v^d,s^d)\le v^\gamma s^{d-\gamma}$.
Here $d-\gamma=19/32$. Dividing by $v^\gamma$ and then taking
the measurable-set supremum proves \eqref{eq:v89-critical-finite}.
Tonelli and monotone exhaustion justify the full sum, since all
terms used in the upper comparison are nonnegative.

Fix $L,s,B_*$ before the collision limsup. Apply
\eqref{eq:v83-fixed-width-input} to the actual fractional source,
and enlarge $B_*$ only afterwards. The capped remainder vanishes,
proving \eqref{eq:v89-critical-ordered}.
\end{proof}

\begin{corollary}[The full surface at the critical exponent]
\label{cor:v89-full-critical}
For fixed $L>0$ and sufficiently small fixed $r>0$,
\begin{align}
 \limsup_m\sup_{R,n,k}\sup_{|J|=L}
       \mathcal N_Q(m^2\overline{\mathcal F}_r;J)
 \le C\{L^{31/32}r^{1/32}+(L+4r)(1+r^{19/32})\}.
                                      \label{eq:v89-full-critical}
\end{align}
In particular the right side is $O(r)$ when $L=r$.
\end{corollary}
\begin{proof}
The surface with $\ell>2r$ is the long-residence source in
\eqref{eq:v88-long-input}; its measurable-set norm is bounded by
$L^{1/Q}$ times its essential height. The complementary surface
is dominated by $Z_{2r}$. Apply Theorem~\ref{thm:v89-endpoint-surface}
and subadditivity. The two classes exhaust the original first-exit
assignment, retaining incident components and backward lifetime.
\end{proof}

The exponent $32/31$ is inherited from the presently available
residence exponent $1/32$, not asserted to be physically sharp.
The new assertion is endpoint boundedness. For a fixed $L$, its
right side need not tend to zero with $r$; no endpoint convergence
or essential-height bound is inferred from it.
